% Living literature review. Edit this file and references.bib together.
% Compile with latexmk -pdf literature_review.tex, or tectonic literature_review.tex.
\documentclass[11pt]{article}
\usepackage[T1]{fontenc}
\usepackage{lmodern}
\usepackage[margin=1in]{geometry}
\usepackage{amsmath,amssymb}
\usepackage{microtype}
\usepackage{needspace}
\usepackage{booktabs,tabularx,array}
\usepackage[numbers,sort&compress]{natbib}
\usepackage{xurl}
\usepackage[colorlinks=true,linkcolor=blue,citecolor=blue,urlcolor=blue]{hyperref}
\setlength{\parindent}{0pt}
\setlength{\parskip}{0.45em}
\setcounter{tocdepth}{1}
\setlength{\emergencystretch}{2em}
\newcommand{\paper}[2]{\Needspace{8\baselineskip}\subsection{#2}\label{paper:#1}\textit{BibTeX key:} \texttt{#1}\quad\citep{#1}\par}
\newcommand{\doi}[1]{\href{https://doi.org/#1}{doi:#1}}
\newcommand{\arxiv}[1]{\href{https://arxiv.org/abs/#1}{arXiv:#1}}
\newcommand{\summary}{\par\textbf{Summary.} }
\newcommand{\relevance}{\par\textbf{Relevance and limits.} }
\title{Literature Review for Coupled Hubbard Ladders}
\author{Ladder MPS and mean field project}
\date{Ladder-material review updated 24 September 2026}

\begin{document}
\maketitle

This review collects 73 papers relevant to an introduction on competing superconducting, charge, and magnetic order in weakly coupled repulsive Hubbard ladders. The 24 September expansion adds a material-by-material comparison, magnetic and electronic parameter benchmarks, pressure and carrier-partition evidence, and the limits of the ideal ladder description. It moves from the general correlated-electron problem to experimental ladder systems, isolated-ladder physics, dimensional crossover, and the closest matrix-product-state plus mean-field (MPS+MF) predecessors. Each annotation explains the result and its use in this project; the bibliography provides complete citations and persistent links.

The project uses a two-leg extended Hubbard ladder with leg hopping $t$, rung hopping $t_0$, onsite repulsion $U$, and nearest-neighbor density interaction $V$, embedded in a transverse mean-field environment. Its principal current regime is $U/t=8$ and density $n=0.9375$, with square and two cubic coupling geometries. Interladder hopping $t_\perp$ generates effective couplings proportional to $t_\perp^2/|E_p|$. These are project definitions, taken from the repository methods notes, rather than findings attributed to the papers below.

\textbf{Reading scope.} This is an introduction-oriented narrative review, with the general search through 7 September 2026 and a focused ladder-cuprate expansion through 24 September, rather than an exhaustive systematic review. Coverage includes the principal copper-oxide ladder families relevant to this Hamiltonian and counterexamples where buckling or interladder coupling changes the physics; molecular, vanadate, and iron-based ladder magnets are outside this scope. Summaries draw on primary abstracts and selected full-text sections, distinguished in the source notes. Three entries remain preprints; the stable key \texttt{Scheie2025} now points to its 2026 journal publication. ``Relevance and limits'' is our synthesis. The historical project-evidence paragraph below retains its own date; current numerical evidence is in the companion manuscript.

\tableofcontents

\section{The Hubbard problem and intertwined electronic orders}
These papers motivate why superconductivity should be studied together with magnetism and charge organization. They supply the broad framing; they do not establish the phase diagram of the present ladder array.

\paper{Hubbard1963}{Hubbard 1963 --- Electron correlations in narrow energy bands}
\doi{10.1098/rspa.1963.0204}
\summary Hubbard introduces a simplified interacting-electron model to describe narrow bands in which local correlations compete with electron motion. The analysis compares a Hartree--Fock treatment with a Green-function approximation that connects atomic and uncorrelated band limits. The paper's central contribution is the model and its account of interaction effects beyond an independent-electron description, including conditions relevant to magnetic behavior. It provides the historical foundation for studying how a small number of microscopic terms can produce qualitatively different electronic states. It is not a theory of superconductivity in doped ladders, and its approximate solution should be distinguished from the Hamiltonian that later work studies with more accurate methods.
\relevance Cite this paper when introducing the Hubbard Hamiltonian. The present $V$ term and transverse reduction are subsequent extensions, so they require their own motivation and citations rather than attribution to the original model.

\paper{Anderson1987}{Anderson 1987 --- The resonating valence bond proposal}
\doi{10.1126/science.235.4793.1196}
\summary Anderson proposes that the proximity of copper-oxide superconductivity to an insulating magnetic state can be understood through resonating valence bonds. In this picture, singlet correlations formed in the insulator provide a starting point for charged pairs when carriers are introduced. The proposal emphasizes electronic and magnetic energy scales, low dimensionality, and the possible role of frustration. Its historical significance is to make pairing from a strongly correlated parent state a concrete alternative to a weakly interacting metal as the starting point. The paper advances a physical mechanism and predictions, rather than a controlled numerical demonstration of superconductivity for all repulsive Hubbard models.
\relevance It supplies a concise historical motivation for repulsion-driven singlet pairing. A spin gap, pair-binding energy, or short-range singlet structure in a ladder remains separate evidence from phase coherence in an array; the RVB proposal does not remove that distinction.

\paper{Lee2006}{Lee, Nagaosa, and Wen 2006 --- Doping a Mott insulator}
\doi{10.1103/RevModPhys.78.17}; \arxiv{cond-mat/0410445}
\summary This review develops the cuprate problem from the perspective of doping a Mott insulator. It discusses spin-charge separation, emergent gauge descriptions, singlet formation, the pseudogap, and the relation of these ideas to $d$-wave superconductivity. A recurring issue is how pairing tendencies and charge coherence develop on different scales when electrons are strongly constrained by interactions. The review gives extensive conceptual and experimental context, while evaluating limitations of explanations based solely on phase fluctuations. Its emphasis is on a particular family of theoretical descriptions of the cuprates, rather than a consensus derivation of their complete phase diagram.
\relevance This is a broad introductory citation for why local pairing, magnetic correlations, and superconducting coherence must be distinguished. Use it selectively: the project's MPS+MF calculation does not itself establish fractionalization, a gauge theory, or a cuprate pseudogap, so those concepts should not be presented as numerical conclusions.

\paper{Scalapino2012}{Scalapino 2012 --- A common thread in unconventional pairing}
\doi{10.1103/RevModPhys.84.1383}; \arxiv{1207.4093}
\summary Scalapino reviews evidence that spin fluctuations provide a common pairing interaction across several families of unconventional superconductors. The discussion connects experimental observations to microscopic model calculations and explains how a predominantly repulsive interaction can favor a superconducting gap that changes sign in momentum space. It places the Hubbard model and numerical pairing studies in a wider materials setting, while keeping the structure of the pairing interaction central. This is especially useful for explaining why strong magnetic correlations can promote pairing even though static magnetic order may compete with the resulting superconducting state. The review does not imply that every antiferromagnetically correlated model is superconducting.
\relevance Use this paper to motivate treating spin and singlet-pair channels together and examining their relative signs. In the ladder calculation, a rung-versus-leg sign change is evidence about a pair form factor; an enhanced pair-field magnitude alone does not identify a pairing mechanism or thermodynamic order.

\paper{Arovas2022}{Arovas and colleagues 2022 --- The Hubbard model}
\doi{10.1146/annurev-conmatphys-031620-102024}; \arxiv{2103.12097}
\summary Arovas, Berg, Kivelson, and Raghu review what is known about the repulsive Hubbard model, emphasizing exact results and controlled approximations in limits with identifiable small parameters. They discuss how ground-state behavior depends on lattice, interaction strength, and filling, and identify substantial unresolved regions despite the Hamiltonian's simplicity. The review is valuable because it distinguishes established statements from plausible extrapolations and competing proposals. Its principal focus is two-dimensional systems, but lower and higher dimensions provide useful comparisons. The message for an introduction is that the Hubbard model is a productive organizing framework whose superconducting and competing-order physics remains sensitive to regime and method.
\relevance This is the strongest broad citation for positioning an anisotropic ladder-array calculation as a tractable part of a harder problem. It also supports stating the limits of the calculation explicitly: a controlled treatment of an individual ladder does not make static transverse mean field exact.

\paper{Tranquada1995}{Tranquada and colleagues 1995 --- Evidence for spin and charge stripes}
\doi{10.1038/375561a0}
\summary Neutron-scattering measurements on a lanthanum-based copper-oxide superconductor reveal signatures interpreted as spatial organization of holes and spins into stripes. The observations connect incommensurate magnetic correlations with charge modulation, making electronic inhomogeneity a central part of the cuprate problem. Rather than treating all doped holes as a spatially uniform background, the stripe picture describes charge-rich regions separating differently phased antiferromagnetic regions. The paper is an experimental landmark for the intertwined charge and magnetic textures that later microscopic calculations try to reproduce. Its material-specific setting, including structural effects that help stabilize order, matters when comparing with ideal lattice models.
\relevance Cite it for the experimental motivation to retain both spin and density channels. It does not justify labeling every density oscillation in an open ladder as a stripe: spatial relations between charge and spin, boundary sensitivity, and persistence with size still have to be examined.

\paper{Fradkin2015}{Fradkin, Kivelson, and Tranquada 2015 --- Theory of intertwined orders}
\doi{10.1103/RevModPhys.87.457}; \arxiv{1407.4480}
\summary This colloquium argues that the relationship among superconductivity, charge order, spin order, and related broken symmetries can involve cooperation and common origins as well as competition. It reviews experimental motivations and theoretical descriptions, including pair-density-wave order and the possibility of new phases emerging from couplings among order parameters. The phrase intertwined orders is useful because coexistence is not automatically a weak perturbation of a single dominant state. Equally, observing several enhanced correlations does not by itself determine which order survives in the thermodynamic limit. The review organizes these possibilities without providing a universal microscopic solution.
\relevance This is a central framing citation for comparing superconducting, CDW, and SDW solutions on equal footing. It motivates allowing coexistence and spatial texture, while the present project's branch energies and correlations must supply the model-specific evidence. It should not be used to assume that every competing branch is part of one equilibrium phase.

\section{Stripe and pairing competition in two dimensional Hubbard calculations}
These studies explain the wider numerical challenge and why small changes of Hamiltonian or texture matter. Cylinders, fully two-dimensional limits, and weakly coupled two-leg ladders are distinct settings.

\paper{Zheng2017}{Zheng and colleagues 2017 --- Stripe order in the underdoped Hubbard model}
\doi{10.1126/science.aam7127}; \arxiv{1701.00054}
\summary This collaborative study uses complementary numerical methods to investigate the underdoped square-lattice Hubbard model in a regime where competing orders are especially difficult to resolve. The calculations favor stripe order and reveal a shallow dependence of energy on stripe wavelength. Changes of stripe period are connected to changes in pairing tendencies, showing that the low-energy competition involves a family of textures rather than a single striped state. The favored filled stripes also differ from those typically inferred in cuprate materials, underscoring the importance of the precise model. Agreement across methods gives this paper particular weight as a benchmark for the numerical difficulty of ordering questions.
\relevance It motivates multiple spatial seeds, matched energy comparisons, and caution about near-degenerate stripe periods. The reported square-lattice result at its studied filling cannot fix the preferred period or phase of an array at $n=0.9375$ with a different transverse Hamiltonian.

\paper{JiangDevereaux2019}{Jiang and Devereaux 2019 --- Pairing, stripes, and next-nearest hopping}
\doi{10.1126/science.aal5304}; \arxiv{1806.01465}
\summary Large-scale DMRG calculations on four-leg Hubbard cylinders at one-eighth hole doping show that next-nearest-neighbor hopping can change the relationship between pairing and charge stripes. For the studied nonzero hoppings, the results are consistent with Luther--Emery behavior: algebraic superconducting and charge correlations coexist with short-range spin correlations. This contrasts with the stronger insulating-stripe character and rapidly decaying pairing found for the nearest-neighbor-only case examined in the paper. The work illustrates how a modest microscopic change can reorganize several competing observables, and why long cylinders and accurate retained-state extrapolations are needed to distinguish those behaviors.
\relevance Cite it as evidence that stripe formation and pairing need not have a universal relationship. Its $t'$ is a diagonal hopping within the square lattice; it is neither the project's rung hopping $t_0$ nor its interladder hopping $t_\perp$. Four-leg cylinder correlations also do not directly establish bulk order in the present array.

\paper{Qin2020}{Qin and colleagues 2020 --- Absence of superconductivity in a studied pure Hubbard regime}
\doi{10.1103/PhysRevX.10.031016}; \arxiv{1910.08931}
\summary Qin and colleagues combine DMRG with constrained-path auxiliary-field quantum Monte Carlo to study pairing in the two-dimensional Hubbard model with nearest-neighbor hopping only. Narrow cylinders provide a setting for detailed cross-method checks, while larger and fully periodic systems help address the thermodynamic limit. In the intermediate-to-strong-coupling regime near optimal hole doping that they investigate, they find a nonsuperconducting ground state. The result is a significant counterweight to assuming that the minimal repulsive model automatically produces robust superconductivity. Its strength comes from systematic comparisons and scaling, and its scope remains the specified Hamiltonian and parameter regime rather than every possible Hubbard model.
\relevance This paper helps motivate explicit tests against insulating competitors and the study of extended interactions or anisotropy. It does not contradict superconductivity in a different ladder array by itself. Avoid paraphrasing the title as a theorem excluding superconductivity at all dopings, couplings, or transverse geometries.

\paper{JiangKivelson2022}{Jiang and Kivelson 2022 --- Stripe order enhanced superconductivity}
\doi{10.1073/pnas.2109406119}; \arxiv{2105.07048}
\summary DMRG calculations on four- and six-leg Hubbard cylinders examine a periodically modulated transverse hopping intended to represent the effects of stripe-like inhomogeneity. Even moderate imposed modulation substantially enhances long-distance superconducting correlations in the studied regimes and produces behavior consistent with a spin-gapped state with superconducting quasi-long-range order. The paper demonstrates that inhomogeneity can strengthen pairing correlations rather than merely disrupt a uniform superconductor. The mechanism is examined in a deliberately modulated Hamiltonian, which makes the relation between texture and pairing experimentally suggestive but also sharply defines what is being demonstrated.
\relevance It provides a useful conceptual bridge to geometry-dependent ladder arrays and a reason to inspect pairing within striped solutions. The modulation is imposed, so the result does not establish spontaneous stripe selection in the project's uniform ladder Hamiltonian. The journal issue is dated 2022, although the article appeared online in December 2021; the bibliography uses the issue year.

\paper{Xu2024}{Xu and colleagues 2024 --- Superconductivity with partially filled stripes}
\doi{10.1126/science.adh7691}; \arxiv{2303.08376}
\summary This study combines DMRG and auxiliary-field quantum Monte Carlo to investigate the Hubbard model including next-nearest-neighbor hopping. It reports superconductivity on both electron- and hole-doped sides, with stronger superconductivity on the hole-doped side coexisting with partially filled stripes. The stripe filling differs from both fully filled stripes in some nearest-neighbor-only calculations and simple half-filled stripes. The analysis emphasizes the exceptionally small energy scales separating competing orders and the need to examine a range of sizes and boundary conditions. It advances the numerical case that superconductivity and nonuniform charge order can coexist in a physically relevant extended hopping model.
\relevance Use it as a modern comparison for coexistence and accurate energetic discrimination. It concerns a different Hamiltonian from both Qin's nearest-neighbor-only model and the present transverse mean-field reduction, so these results should be presented as regime-dependent findings rather than an apparent contradiction. The journal publication is 2024; the preprint first appeared in 2023.

\paper{Agterberg2020}{Agterberg and colleagues 2020 --- The physics of pair-density waves}
\doi{10.1146/annurev-conmatphys-031119-050711}; \arxiv{1904.09687}
\summary This review describes superconducting states whose pair order parameter varies in space at a nonzero center-of-mass wave vector. It discusses their symmetry properties, the charge and other orders they can induce, possible vestigial phases, experimental evidence, and microscopic proposals. It also distinguishes the cuprate-related pair-density-wave problem from other settings with modulated pairing, including weak-coupling finite-momentum states. The review is valuable for deciding what an observed pair-field texture actually means: spatially varying magnitude, internal orbital sign structure, and modulation of the superconducting phase are related but different pieces of information.
\relevance It supplies terminology for the project's sign-resolved pair fields and Fourier analysis. A rung-versus-leg $d$-wave sign change concerns the internal pair form factor, while a PDW concerns center-of-mass modulation. Neither a seeded node nor a nonzero Fourier component in a finite open system is sufficient to establish a thermodynamic PDW phase.

\section{Ladder materials and experimental motivation for extended interactions}
\label{sec:materials-review}
The material comparison must separate four questions: whether the structure
contains ladders, whether its spins realize an approximately isolated ladder,
where doped holes reside, and whether those carriers establish superconducting
coherence. The same chemical family can provide different benchmarks for each.
Table~\ref{tab:material-map} summarizes this distinction; the annotations below
provide the primary evidence and its limits.

\begin{table}[htbp]
\centering\small
\caption{Core ladder-cuprate families and their role in the project. A magnetic
gap is an excitation energy, not a superconducting transition temperature.}
\label{tab:material-map}
\begin{tabularx}{\linewidth}{@{}p{0.28\linewidth}X@{}}
\toprule
Material & Evidence and model relevance \\
\midrule
SrCu$_2$O$_3$; Sr$_2$Cu$_3$O$_5$ & Two- versus three-leg insulating benchmarks: gapped versus gapless spin response in the original study~\citep{Azuma1994}. \\
La$_6$Ca$_8$Cu$_{24}$O$_{41}$ & Nominally hole-free chain--ladder reference; Raman constrains anisotropic exchange and cyclic exchange~\citep{Nucker2000,Schmidt2005}. \\
La$_4$Sr$_{10}$Cu$_{24}$O$_{41}$ & Approximately undoped ladders despite residual chain holes; resolved one-/two-triplon neutron spectra~\citep{Notbohm2007}. \\
Sr$_{14-x}$Ca$_x$Cu$_{24}$O$_{41}$ & Intrinsically hole-bearing chains and ladders; charge order, carrier redistribution, and pressure-induced superconductivity~\citep{Uehara1996,Abbamonte2004,Nucker2000}. \\
La$_{1-x}$Sr$_x$CuO$_{2.5}$ & Doped coupled ladders; metallicity without superconductivity down to 5 K in the cited experiment~\citep{Hiroi1995,NormandRice1996}. \\
CaCu$_2$O$_3$ & Buckled ladder-like structure; weakened rung path and magnetic order at 25 K, not a standard spin-gap ladder~\citep{Kiryukhin2001}. \\
\bottomrule
\end{tabularx}
\end{table}

\subsection*{Insulating ladders, exchange anisotropy, and structural limits}

\paper{Azuma1994}{Azuma and colleagues 1994 --- Two- and three-leg magnetic benchmarks}
\doi{10.1103/PhysRevLett.73.3463}
\summary Susceptibility and copper NMR compare SrCu$_2$O$_3$ with Sr$_2$Cu$_3$O$_5$. The two-leg material has an activated response, with a susceptibility-derived gap of about 420 K; the three-leg material has a gapless response. These observations establish the even/odd ladder distinction in copper oxides.
\relevance This is evidence for insulating spin physics, not a superconducting transition at 420 K or a measurement of hole binding. The quoted gap is tied to the probe and analysis; it should not be imposed as an exact parameter constraint on every two-leg material.

\paper{Eccleston1998}{Eccleston and colleagues 1998 --- Distinguishing ladder and chain excitations}
\doi{10.1103/PhysRevLett.81.1702}
\summary Neutron spectroscopy resolves the higher-energy ladder response and lower-energy dimer-chain excitations of Sr$_{14}$Cu$_{24}$O$_{41}$. The original ladder dispersion analysis gave a gap near 32.5 meV and unequal leg and rung exchange, while the roughly 11.5 meV feature belonged to chain dimers.
\relevance This establishes why the two subsystems cannot be conflated. Early exchange estimates from this analysis are not interchangeable with later fits that retain cyclic exchange, nor with the distinct La-substituted sample below.

\paper{Nunner2002}{Nunner and colleagues 2002 --- Cyclic exchange in optical spectra}
\doi{10.1103/PhysRevB.66.180404}; \arxiv{cond-mat/0203472}
\summary Dynamical DMRG analyzes how four-spin exchange changes one-triplet and singlet two-triplet excitations. For the studied (La,Ca)$_{14}$Cu$_{24}$O$_{41}$ spectra, simultaneous agreement with the spin gap and optical bound-state spectrum requires a substantial cyclic term and unequal leg/rung exchange.
\relevance A spin gap alone does not determine a unique nearest-neighbor exchange ratio. The paper writes cyclic exchange using a permutation operator, which also generates bilinear terms; coefficients must be converted before comparison with a convention retaining only explicit four-spin products.

\paper{Schmidt2005}{Schmidt and colleagues 2005 --- Raman comparisons across three ladder materials}
\doi{10.1103/PhysRevB.72.094419}; \arxiv{cond-mat/0503558}
\summary A common Heisenberg-plus-four-spin analysis compares SrCu$_2$O$_3$, Sr$_{14}$Cu$_{24}$O$_{41}$, and La$_6$Ca$_8$Cu$_{24}$O$_{41}$. The fitted $J_{\rm leg}/J_{\rm rung}$ is 1.5 for the first two and 1.2 for the last, with $J_{\rm cyc}/J_{\rm rung}=0.2$ in that convention. Exchange anisotropy helps explain the differing Raman peak widths.
\relevance La$_6$Ca$_8$ and La$_4$Sr$_{10}$ are distinct references, not alternate names for one sample. These fits support stronger-leg parameter regimes but do not supply a universal hopping ratio or a pressure-dependent superconducting Hamiltonian.

\paper{Notbohm2007}{Notbohm and colleagues 2007 --- The La4Sr10 neutron benchmark}
\doi{10.1103/PhysRevLett.98.027403}; \arxiv{cond-mat/0608109}
\summary Opposite rung-parity modulations separate one- and two-triplon neutron scattering in La$_4$Sr$_{10}$Cu$_{24}$O$_{41}$, whose residual holes predominantly occupy chains. The ladder gap is $26.4\pm0.3$ meV. A fit gives $J_{\rm leg}=186$, $J_{\rm rung}=124$, and $J_{\rm cyc}=31$ meV, using the paper's explicit four-spin convention. The authors also give approximate Hubbard values $t_{\rm leg}=0.42$, $t_{\rm rung}=0.34$, and $U=3.72$ eV.
\relevance The reported hopping estimate gives $t_0/t\simeq0.81$; a leading-order square-root conversion of the exchanges gives 0.82. Both are model-dependent estimates, not direct hopping measurements. Stronger leg coupling occurs despite a longer Cu--Cu leg distance (3.96 versus 3.84 \AA), illustrating the role of the oxygen environment. This insulating benchmark does not calibrate a superconducting Ca-rich sample.

\paper{Bordas2006}{Bordas and colleagues 2006 --- Ab initio electronic parameters}
\doi{10.1007/s00214-006-0099-z}
\summary Embedded-cluster configuration-interaction calculations determine exchange and hopping parameters for SrCu$_2$O$_3$, Sr$_2$Cu$_3$O$_5$, and CaCu$_2$O$_3$, including longer-range and cyclic processes. Direct exchange, charge-transfer hopping, and onsite repulsion are examined alongside effective magnetic couplings.
\relevance This supplies a microscopic complement to spectroscopy and warns against deriving every coupling from bond length alone. Its neutral--ionic hopping symbol $t_0$ is not our rung-hopping definition. We cite the approach and compound coverage, without transferring unverified numerical table entries.

\paper{Kiryukhin2001}{Kiryukhin and colleagues 2001 --- CaCu2O3 as a structural counterexample}
\doi{10.1103/PhysRevB.63.144418}; \arxiv{cond-mat/0009158}
\summary Susceptibility and neutron diffraction find chain-like magnetism and an incommensurate ordered state below $T_N=25$ K in CaCu$_2$O$_3$. The rung Cu--O--Cu angle is about $123^\circ$, strongly departing from the nearly straight path in the standard ladder compounds.
\relevance The structural label ``ladder'' does not guarantee a dominant-rung singlet state or even a good isolated two-leg description. $T_N$ is magnetic, not superconducting; buckling and competing interladder interactions need explicit treatment for this compound.

\paper{Hiroi1995}{Hiroi and Takano 1995 --- Doping does not guarantee superconductivity}
\doi{10.1038/377041a0}
\summary Sr substitution in LaCuO$_{2.5}$ introduces holes and an insulator-to-metal transition, but the reported samples do not become superconducting down to 5 K.
\relevance This is a useful negative material result. It does not exclude pairing in all ladder Hamiltonians, nor justify describing every hole-doped ladder cuprate as a superconductor.

\paper{NormandRice1996}{Normand and Rice 1996 --- Three-dimensional coupling in LaCuO2.5}
\doi{10.1103/PhysRevB.54.7180}; \arxiv{cond-mat/9603054}
\summary Electronic-structure and magnetic modeling find substantial coupling between the ladder units of LaCuO$_{2.5}$. Its three-dimensional environment is central to interpreting the material.
\relevance The result motivates varying transverse connectivity, while preventing a literal identification of this material with our schematic cubic kernels. A geometrical ladder motif is not sufficient evidence for a weak-coupling reduction.

\subsection*{Carrier partition, superconductivity, and competing order}

\paper{Nucker2000}{N\"ucker and colleagues 2000 --- Where the holes reside}
\doi{10.1103/PhysRevB.62.14384}; \arxiv{cond-mat/0008320}
\summary Polarization-dependent x-ray absorption distinguishes chain and ladder oxygen-hole states under Sr/Ca and La/Y substitutions. Isovalent Ca substitution preserves the nominal total hole count while changing orbital occupation; trivalent substitution reduces that count. Partial ladder occupation occurs in Sr$_{14}$Cu$_{24}$O$_{41}$, whereas residual holes in less hole-rich members can reside predominantly on chains.
\relevance At O$_{41}$ stoichiometry the all-divalent family has six formal holes per formula unit, La$_4$Sr$_{10}$ has two, and La$_6$Ca$_8$ has none, by charge neutrality relative to Cu$^{2+}$. These are total counts, not ladder densities. Oxygen nonstoichiometry and probe-dependent carrier partition prevent an automatic conversion from composition to $n$.

\paper{Uehara1996}{Uehara and colleagues 1996 --- Superconductivity in a ladder material}
\doi{10.1143/JPSJ.65.2764}
\summary Uehara and colleagues report pressure-induced superconductivity in the calcium-rich ladder material $\mathrm{Sr}_{0.4}\mathrm{Ca}_{13.6}\mathrm{Cu}_{24}\mathrm{O}_{41.84}$. Electrical measurements show superconducting transitions under applied pressure, and magnetization provides additional evidence for a superconducting fraction. This landmark observation gives experimental substance to the proposal that doped spin ladders can be relevant building blocks for superconductivity. The material is more complicated than a single isolated two-leg Hubbard ladder: its structural units, carrier distribution, and pressure response all matter. The experiment therefore establishes the relevance of the ladder-material family without uniquely identifying a minimal microscopic pairing Hamiltonian.
\relevance It is an appropriate introductory citation for studying coupled ladders as more than a numerical simplification. It does not calibrate the project's $t_0$, $V$, or $t_\perp$, and the existence of a superconducting transition in a crystal cannot be transferred directly to one finite isolated ladder.

\paper{Nagata1998}{Nagata and colleagues 1998 --- Pressure-induced dimensional crossover}
\doi{10.1103/PhysRevLett.81.1090}
\summary Anisotropic resistivity measurements under pressure investigate how superconductivity emerges in a hole-doped two-leg ladder compound. Superconductivity is accompanied by metallic transport between ladders, in contrast to the more insulating transverse response at ambient pressure. The authors interpret this as evidence that pressure promotes a dimensional crossover and connects a low-dimensional insulating state to a superconducting state in an anisotropic higher-dimensional system. The paper is especially relevant because it measures a property associated with coupling between ladders rather than only signatures inside an individual ladder. Its interpretation remains tied to the pressure-dependent material, where several microscopic parameters may change simultaneously.
\relevance This is a direct experimental motivation for asking how interladder coupling changes the ordering of a paired subsystem. Pressure should not be represented as a change in only one hopping unless supported by a material model. The present geometry comparisons are theoretical changes of Hamiltonian, not a quantitative reconstruction of this pressure trajectory.

\paper{Radheep2013}{Mohan Radheep and colleagues 2013 --- Uniaxial-pressure onset report}
\arxiv{1303.0921}\quad\textit{Preprint}
\summary The authors report a resistive onset up to 24 K in Sr$_3$Ca$_{11}$Cu$_{24}$O$_{41}$ under approximately 0.06 GPa compression along $c$. Their magnetization example instead has an onset near 6.8 K under approximately 0.05 GPa along $a$, with an estimated superconducting fraction near 7\% at 5 K.
\relevance This pressure-direction-dependent report must be distinguished from the 12 K, 3 GPa benchmark of Uehara et al. The cited magnetic measurement is not bulk confirmation at 24 K under the same conditions. No journal reference appears on the checked arXiv record; we retain its preprint status and resistive-onset qualification.

\paper{Fujiwara2003}{Fujiwara and colleagues 2003 --- Superconducting-state NMR}
\doi{10.1103/PhysRevLett.90.137001}; \arxiv{cond-mat/0303278}
\summary High-pressure NMR on Sr$_2$Ca$_{12}$Cu$_{24}$O$_{41}$ examines the superconducting quasiparticle response and its relationship to spin-gap formation. The reported gapped behavior is described as $s$-wave-like.
\relevance The result is not a phase-sensitive measurement of leg-versus-rung pair signs. A ladder state with opposite bond signs need not have nodes on its actual Fermi surfaces; a spin gap, a quasiparticle gap, and a pair-breaking energy are separate observables.

\paper{Fujiwara2009}{Fujiwara and colleagues 2009 --- Pressure-dependent charge transfer and spin response}
\doi{10.1103/PhysRevB.80.100503}; \arxiv{0908.2845}
\summary NMR and NQR follow Sr$_2$Ca$_{12}$Cu$_{24}$O$_{41}$ up to the optimal pressure near 3.8 GPa. The quadrupole response supports pressure-induced hole transfer from chains to ladders, while the normal-state magnetic response also evolves.
\relevance Pressure changes ladder filling as well as coupling. A fixed-density transverse-hopping scan isolates one mechanism but does not reproduce the complete experimental pressure trajectory.

\paper{Frank2014}{Frank and colleagues 2014 --- Optical charge dynamics under pressure}
\doi{10.1103/PhysRevB.90.224516}
\summary Polarized infrared reflectivity on Sr$_{2.5}$Ca$_{11.5}$Cu$_{24}$O$_{41}$ investigates pressure-dependent charge dynamics, redistribution between chains and ladders, and anisotropy. Absolute ladder occupations depend on the experimental interpretation and calibration.
\relevance Optical spectral weight supports pressure-dependent mobility and carrier transfer, but does not uniquely identify the model density $n=15/16$. Chemical and applied pressure should not be treated as interchangeable changes of a single parameter.

\paper{Vuletic2003}{Vuleti\'c and colleagues 2003 --- Suppression of a CDW scale by Ca substitution}
\doi{10.1103/PhysRevLett.90.257002}; \arxiv{cond-mat/0305159}
\summary Transport, dielectric, and optical measurements find a CDW transition scale falling from about 210 K at $x=0$ to 10 K at $x=9$ in Sr$_{14-x}$Ca$_x$Cu$_{24}$O$_{41}$. This is a charge-order scale, not a superconducting $T_c$.
\relevance The trend motivates competing charge and pairing channels, but does not establish the absence of charge order at all larger Ca concentrations; resonant scattering supplies a complementary test.

\paper{Abbamonte2004}{Abbamonte and colleagues 2004 --- Crystallization of charge holes in a spin ladder}
\doi{10.1038/nature02925}; \arxiv{cond-mat/0501087}
\summary Resonant x-ray scattering identifies charge organization associated with the ladder subsystem of $\mathrm{Sr}_{14}\mathrm{Cu}_{24}\mathrm{O}_{41}$. The results are interpreted as a hole crystal, providing experimental evidence for an electronically ordered state in a ladder material that had also motivated theories of pairing. The work connects low-dimensional correlation physics to directly observable charge modulation and establishes that the fate of doped carriers is not limited to uniform transport or superconductivity. Its experimental sensitivity to electronic contrast is central to the interpretation, while the detailed ordering pattern belongs to a specific compound and carrier distribution.
\relevance This paper motivates retaining CDW competitors when studying superconducting ladder arrays. It should be paired with boundary-scaling references when discussing numerical density profiles: an experimentally detected bulk charge pattern and open-boundary Friedel oscillations are different kinds of evidence. The project's density and interaction values should not be inferred from the compound's chemical formula alone.

\paper{Rusydi2006}{Rusydi and colleagues 2006 --- Commensurability and hole crystals}
\doi{10.1103/PhysRevLett.97.016403}; \arxiv{cond-mat/0511524}
\summary Resonant scattering detects commensurate ladder hole crystals at selected compositions, including $x=10,11,12$, and relates their appearance to commensurability. Longitudinal periods three and five occur in the studied family.
\relevance Ca content alone does not determine whether a CDW exists. Longitudinal charge modulation and coherence between ladders do not establish unequal average fillings on different ladders, the channel excluded by the current equal-density ansatz.

\paper{Katano1999}{Katano and colleagues 1999 --- Pressure and static magnetic order}
\doi{10.1016/S0921-4526(98)00692-9}
\summary Neutron measurements on Sr$_{2.5}$Ca$_{11.5}$Cu$_{24}$O$_{41}$ find a pressure-suppressed ordered magnetic moment, apparently disappearing near the superconducting threshold.
\relevance This supports separating static magnetism from the dynamic spin gap. It does not demonstrate universal mutual exclusion of magnetic and superconducting order across the whole family or directly validate a particular mean-field texture.

\subsection*{Electronic models, localization, and transverse geometry}

\paper{Popovic2000}{Popovi\'c and colleagues 2000 --- An isotropic optical-model benchmark}
\doi{10.1103/PhysRevB.62.4963}; \arxiv{cond-mat/0005096}
\summary An optical model of Sr$_{14}$Cu$_{24}$O$_{41}$ uses $t=t_0=0.26$ eV, $U=2.1$ eV, and interladder hopping $t_{xy}=0.026$ eV. The latter enters the two displaced paths between neighboring ladders.
\relevance This motivates the project's isotropic-rung trellis comparison with two path amplitudes of $0.1t$. It is a different effective parameterization from the neutron and RIXS models; it does not independently determine $V$ or calibrate a high-pressure superconducting sample.

\paper{Gelle2006}{Gell\'e and Lepetit 2006 --- Incommensurate structural modulation}
\doi{10.1103/PhysRevB.74.235115}; \arxiv{cond-mat/0411135}
\summary Correlated ab initio calculations derive spatially varying electronic parameters for the ladder subsystem of the incommensurate Sr/Ca compounds. Strong onsite-energy modulation along the legs is found in the highly Ca-substituted case, while the two sites of a rung remain approximately equivalent in the reported calculation.
\relevance A uniform ladder Hamiltonian omits a potential localization mechanism. Weak direct chain--ladder hybridization does not prevent the chains from setting the electrostatic environment or redistributing charge. The arXiv record contains a mismatched related DOI; the bibliography uses the verified PRB publication above.

\paper{Tseng2022}{Tseng and colleagues 2022 --- Doping, disorder, and magnetic localization}
\doi{10.1038/s41535-022-00502-1}
\summary RIXS compares the high-energy magnetic response of Sr$_{14-x}$Ca$_x$Cu$_{24}$O$_{41}$ across Ca substitution. The interpretation invokes off-ladder impurity potentials and hole localization to account for changes from collective triplon behavior.
\relevance This provides a disorder-based comparison for later attractive-interaction models. A material spectrum can reflect localization as well as intrinsic pairing; neither mechanism is resolved by a uniform static anomalous amplitude alone.

\paper{Hirthe2023}{Hirthe and colleagues 2023 --- Magnetically mediated hole pairing with ultracold atoms}
\doi{10.1038/s41586-022-05437-y}; \arxiv{2203.10027}
\summary A quantum gas microscope realizes doped antiferromagnetic ladders with mixed-dimensional couplings and observes hole pairing associated with magnetic correlations. Engineering the coupling pattern suppresses short-distance Pauli-blocking effects and enhances the binding energy, making compact pairs on the same rung visible in real space. At higher doping, spatial correlations between pairs indicate that the resulting composite objects also interact. The experiment supplies unusually direct evidence that mobile holes can bind in a magnetically correlated background. It provides a controlled realization of an important ingredient in unconventional superconductivity, rather than a measurement of a bulk superconducting transition.
\relevance The paper motivates calculating pair-binding energies and examining pair structure independently of long-range order. Its mixed-dimensional ladder has suppressed rung particle tunneling while retaining exchange, so it differs from an ordinary Hubbard ladder with finite $t_0$. It is especially relevant when comparing the project with the 2026 MPS+MF quantum-simulator preprint.

\paper{Chen2021}{Chen and colleagues 2021 --- Strong near-neighbor attraction in cuprate chains}
\doi{10.1126/science.abf5174}; \arxiv{2106.14272}
\summary This experiment synthesizes and studies doped one-dimensional cuprate chains over a range of carrier concentrations. Angle-resolved photoemission resolves the evolution of spinon and holon features and reveals a prominent spectral branch that the simple onsite Hubbard model does not reproduce quantitatively. Including an additional strong nearest-neighbor attraction explains the observed spectral behavior across the accessible dopings. Coupling to phonons is discussed as a possible origin of that effective attraction. The result is significant because it gives an experimental reason to examine an extended Hubbard Hamiltonian instead of treating its nonlocal interaction solely as a theoretical tuning parameter.
\relevance It motivates the negative-$V$ part of the project's survey. The inferred attraction is an effective interaction in a chain and does not directly determine its value or bond dependence in a ladder. Neither the photoemission observation nor the added attraction alone establishes singlet superconductivity in the present model.

\paper{Padma2025}{Padma and colleagues 2025 --- Beyond-Hubbard pairing in a cuprate ladder}
\doi{10.1103/PhysRevX.15.021049}; \arxiv{2501.10287}
\summary Resonant inelastic x-ray scattering and DMRG examine magnetic excitations in $\mathrm{Sr}_{14}\mathrm{Cu}_{24}\mathrm{O}_{41}$. The measured magnetic response of doped carriers is strongly suppressed relative to the onsite Hubbard model. Including a nearest-neighbor attraction of roughly $V/t=-1$ to $-1.25$ improves agreement with the spectra. The authors attribute this suppression to tighter hole pairing, with substantially increased binding energy and retained $d$-wave-like rung-versus-leg sign structure. Their $V=0$ doped-ladder baseline uses $U/t=8$ and hole density $1/16$; the attraction scan instead adjusts $U$ to maintain $U-V=8t$. The study provides a direct spectroscopic motivation for an extended ladder Hamiltonian, while its main evidence concerns bound carriers rather than a superconducting transition.
\relevance This is a high-priority citation for the $V$ scan. Its fitted rung hopping is $0.84t$, and its main-text model includes diagonal hopping $t'=-0.3t$. The project lacks that diagonal term and holds $U/t=8$ as $V$ varies, so the scans follow different paths in parameter space. The material fit does not validate the project's transverse treatment; stronger binding also need not imply greater pair mobility or a higher ordering temperature.

\paper{Scheie2025}{Scheie and colleagues 2026 --- Cooper-pair localization in ladder magnetic dynamics}
\doi{10.1103/1s7m-xkby}; \arxiv{2501.10296}\quad\textit{Original 2025 key retained}
\summary High-resolution neutron spectroscopy and DMRG investigate spin dynamics in $\mathrm{Sr}_{2.5}\mathrm{Ca}_{11.5}\mathrm{Cu}_{24}\mathrm{O}_{41}$. The spectra lack the incommensurate magnetic signatures expected from mobile doped holes. The proposed explanation combines an attractive nearest-neighbor interaction of order $-t$ with pair localization at pinning centers. In the finite-cluster calculation, open boundaries emulate such pinning and accumulate pairs near the ends, leaving a lower effective hole density in the bulk. This restores a predominantly commensurate magnetic response. The study complements the related x-ray work, but localization is essential to its interpretation: binding alone does not explain every observed feature or establish coherent transport.
\relevance This is particularly relevant to boundary-sensitive density profiles and the distinction between pairing and superconductivity. The author manuscript uses $t=0.4$ eV, $t_0/t=0.84$, and diagonal $t'/t=-0.3$. Its localization mechanism is explicit, so the result is not a homogeneous bulk phase diagram. The publisher record is now Phys. Rev. B \textbf{114}, 144511 (23 September 2026); detailed parameter statements here refer to the checked author manuscript.

\paper{SchmidtUhrig2007}{Schmidt and Uhrig 2007 --- Magnetic signatures of trellis coupling}
\doi{10.1103/PhysRevB.75.224414}; \arxiv{cond-mat/0701551}
\summary A spin-model treatment computes the two-dimensional one-triplon dispersion and neutron structure factor for frustrated trellis arrays relevant to SrCu$_2$O$_3$ and the chain--ladder family. It identifies momentum-dependent signatures of interladder exchange.
\relevance Weak or frustrated transverse coupling does not mean it is absent. This is a magnetic-excitation benchmark, not a determination of electronic pair transfer or validation of the project's static transverse approximation.

\paper{Faye2015}{Faye and colleagues 2015 --- Electronic pairing on the trellis}
\doi{10.1103/PhysRevB.91.195126}
\summary Variational cluster approximation and cellular dynamical mean field theory study a repulsive Hubbard model on coupled two-leg ladders. Both find a superconducting doping dome in the examined regime. A complex, time-reversal-breaking solution appears in the variational-cluster treatment but not in the cellular-DMFT solution.
\relevance This is a direct electronic trellis comparator. Method-dependent phase structure, small clusters, and the different Hamiltonian prevent quantitative transfer of its boundaries. Our real-field ansatz cannot test its complex order parameter; its $t'$ labels interladder hopping, unlike the diagonal $t'$ in the RIXS fits.

\paper{Sushkov2008}{Sushkov and colleagues 2008 --- Orthorhombicity is not a universal hopping prescription}
\doi{10.1103/PhysRevB.77.035124}; \arxiv{0710.5325}
\summary For insulating orthorhombic La$_{2-x}$Sr$_x$CuO$_4$, electronic-structure calculations find a difference of about 8 meV between hoppings along the two plaquette diagonals. The orthorhombic crystallographic directions in this example are diagonal to the nearest-neighbor Cu network.
\relevance This planar comparison prevents equating every orthorhombic distortion with one nearest-neighbor anisotropy. A square-connectivity approximation and isotropic hopping are separate assumptions; ladder rung/leg inequivalence also exists without a small planar orthorhombic perturbation.

\paper{Gao1994}{Gao and colleagues 1994 --- A planar-cuprate transition-temperature comparison}
\doi{10.1103/PhysRevB.50.4260}
\summary Resistive measurements on optimally doped Hg-based planar cuprates find transition temperatures up to 164 K under quasihydrostatic pressure. These compounds provide a high-temperature comparison to the roughly 10 K hydrostatic-pressure ladder benchmark.
\relevance The comparison concerns observed transition temperatures under different pressures, not equal microscopic couplings or a ratio of pair-binding energies. Pressure, sample composition, and onset versus bulk criteria must accompany quoted $T_c$ values.

\subsection*{What the material review implies for the present parameter scan}
The three useful benchmarks are a neutron-derived stronger-leg estimate near
$t_0/t=0.8$ in La$_4$Sr$_{10}$, the explicitly chosen ratio 0.84 in the
Sr/Ca spectroscopic Hubbard models, and the isotropic optical model with
$t_0/t=1$~\citep{Notbohm2007,Padma2025,Scheie2025,Popovic2000}.
They must remain separate parameterizations, including their $U$, diagonal
hopping, filling, cyclic exchange, and disorder assumptions. In the leading
half-filled one-band expansion with a common denominator,
\begin{equation}
 J_\mu\simeq\frac{4t_\mu^2}{U-V_\mu},\qquad
 \frac{t_0}{t}\simeq\sqrt{\frac{J_{\rm rung}}{J_{\rm leg}}}
 \quad(V_{\rm rung}=V_{\rm leg}).
\end{equation}
This is our conditional conversion, not a direct experimental extraction.
Finite-$U$ Hubbard calculations already generate higher-order exchange
processes; adding a spin-model ring term independently can double count them.
The existing $t_0/t=1.2$--1.4 scans explore stronger-rung physics and are not
material fits justified by orthorhombicity. At 1.4 the leading exchange ratio
would be 1.96, opposite to the cited stronger-leg benchmarks.

Neither the magnetic spin gap nor a pairing scale is $T_c$. Uehara's reported
onsets are 12 K at 3 GPa and 9 K at 4.5 GPa; the 24 K uniaxial onset has the
separate qualifications above. An isolated short-range ladder has no
finite-temperature superconducting transition. Transverse coherence,
localization, and charge order therefore matter even when local pairing is
strong. No universal pressure-to-$t_0/t$, pressure-to-$t_\perp$, or
composition-to-density map follows from this review.

\section{Low dimensional constraints and the origin of ladder pairing}
The two-leg ladder is useful because it can retain strong local correlations while supporting a spin-gapped, conducting charge sector. This group establishes that starting point and its limits.

\paper{MerminWagner1966}{Mermin and Wagner 1966 --- Dimensional restrictions on magnetic order}
\doi{10.1103/PhysRevLett.17.1133}
\summary Mermin and Wagner prove the absence of ferromagnetic or antiferromagnetic long-range order at nonzero temperature in one- and two-dimensional isotropic Heisenberg models with suitable short-range interactions. The result demonstrates that fluctuations and dimensionality can invalidate a conventional mean-field ordering picture even when local correlations are strong. Its importance extends beyond a particular material: assumptions about symmetry, interaction range, and temperature determine what kind of ordered state is possible. The original theorem concerns the specified spin models and finite-temperature limit; it should not be quoted as a blanket prohibition of every ordered state in low dimensions.
\relevance This is a useful constraint citation if the introduction contrasts isolated subsystems with higher-dimensional arrays. The current project compares zero-temperature states, so the theorem does not rule out its two-dimensional ground-state order. It also does not by itself prove that a finite ladder is a Luther--Emery liquid or validate a mean-field critical temperature.

\paper{LutherEmery1974}{Luther and Emery 1974 --- Backward scattering in the one-dimensional electron gas}
\doi{10.1103/PhysRevLett.33.589}
\summary Luther and Emery solve a one-dimensional interacting-electron problem at a special attractive backward-scattering coupling and use scaling arguments to describe the broader attractive-backscattering regime. The work underlies the concept of a Luther--Emery liquid, in which the spin sector is gapped while a charge mode remains gapless. This separation provides a framework for understanding how strong pairing correlations can coexist with the absence of a conventional broken-symmetry superconducting state in an isolated one-dimensional system. Density and superconducting correlations then compete through their long-distance behavior rather than through a simple finite order parameter.
\relevance This is the foundational citation for the low-energy language used in ladder studies. Repulsive two-leg Hubbard ladders can realize an analogous low-energy structure, but that is a result of subsequent ladder theory and numerics, not the original continuum solution. Gap measurements and compatible charge and pair exponents are needed to identify that regime in the project's bare ladders.

\paper{Dagotto1992}{Dagotto, Riera, and Scalapino 1992 --- Superconductivity in ladders and coupled planes}
\doi{10.1103/PhysRevB.45.5744}
\summary This early numerical and theoretical study examines two coupled chains or planes described by a $t$--$J$-type Hamiltonian. For substantial coupling between the subsystems, it finds a spin gap, bound holes after doping, and superconducting pairing correlations on finite clusters. The relevant pair is a spin singlet with carriers distributed across the coupled units. The paper establishes a concrete mechanism by which a ladder geometry and magnetic exchange can make pairing favorable even when double occupancy is strongly restricted. It helped turn ladders into a central testing ground for magnetically mediated pairing.
\relevance Cite it for the historical connection among a ladder spin gap, hole binding, and singlet pairing. The $t$--$J$ model and finite-cluster evidence are not identical to the finite-$U$ extended Hubbard model used here. In particular, a bound two-hole state does not by itself determine which ordered many-particle phase wins after ladders are coupled.

\paper{DagottoRice1996}{Dagotto and Rice 1996 --- Ladders between one and two dimensions}
\doi{10.1126/science.271.5249.618}; \arxiv{cond-mat/9509181}
\summary This review explains why assembling spin chains into ladders does not produce a smooth interpolation to a two-dimensional magnet. Even-leg and odd-leg ladders can have qualitatively different magnetic spectra, with even-leg ladders supporting a finite spin gap and short-range magnetic correlations. It connects theoretical predictions to the emerging ladder-material experiments and discusses why doped holes in such structures may bind. The two-leg ladder is therefore presented as a physically distinctive system, not merely the narrowest convenient numerical approximation to a plane. That perspective remains central to interpreting modern Hubbard-ladder calculations.
\relevance This is an efficient broad-to-narrow transition citation for an introduction. It motivates choosing two-leg building blocks while warning against a naive extrapolation from their behavior to arbitrary widths or geometries. Material realizations, isolated-ladder pairing correlations, and superconductivity in a coupled array should be kept as separate stages of the argument.

\paper{Noack1994}{Noack, White, and Scalapino 1994 --- Correlations in a two-chain Hubbard model}
\doi{10.1103/PhysRevLett.73.882}; \arxiv{cond-mat/9401013}
\summary Early DMRG calculations investigate spin and pair-field correlations in the two-chain Hubbard model. At half filling, both types of correlation decay exponentially, consistent with a gapped insulating state. Upon hole doping, magnetic correlations become incommensurate while a spin gap persists, and the pair correlations appear compatible with power-law behavior. The work establishes an important numerical foundation for the doped repulsive ladder as a spin-gapped conductor. It also illustrates the historical difficulty of determining pairing dominance: the apparent pair decay in the accessible systems did not yet provide the long-distance and truncation control available in later studies.
\relevance Cite it as an original numerical predecessor, then use Dolfi and Shen for the modern exponent analysis. The presence of a spin gap and algebraic pairing is already a substantial result, but it neither guarantees pairing dominates density correlations nor proves superconducting long-range order in an isolated finite-width ladder.

\paper{BalentsFisher1996}{Balents and Fisher 1996 --- Weak-coupling phases of the two-chain Hubbard model}
\doi{10.1103/PhysRevB.53.12133}; \arxiv{cond-mat/9503045}
\summary Balents and Fisher develop a controlled weak-coupling analysis for a finite number of coupled interacting electron chains and work out the two-chain Hubbard phase diagram in detail. The classification uses the number of gapless charge and spin modes, providing the familiar language of phases such as C1S0. Interchain hopping and doping alter which scattering processes are important and which sectors become gapped. The analysis identifies circumstances in which a two-chain system behaves as a one-dimensional analogue of a superconductor, while also revealing several competing low-energy regimes.
\relevance This is a primary theoretical citation for the expected Luther--Emery structure and for the importance of rung hopping. The weak-coupling control concerns the intraladder interaction, whereas the project uses $U/t=8$ and a separate weak-interladder approximation. Those are different expansions; the paper's boundaries cannot be used as quantitative phase boundaries for the current strong-coupling survey.

\paper{Noack1997}{Noack and colleagues 1997 --- Enhanced d-wave pairing in the two-leg Hubbard ladder}
\doi{10.1103/PhysRevB.56.7162}; \arxiv{cond-mat/9612165}
\summary DMRG and quantum Monte Carlo calculations examine how pairing correlations depend on the ratio of rung to leg hopping, onsite interaction, and filling. Pairing is enhanced in a regime where the relevant bonding and antibonding band features lie close to the Fermi level. Single-particle spectral weight and antiferromagnetic correlations help interpret this behavior, connecting a change in ladder geometry to the structure of the pairing response. The work provides a microscopic reason to investigate anisotropic ladders rather than assume the isotropic point is universally optimal. Its $d$-wave-like characterization concerns the relative structure of pairing amplitudes on different bonds.
\relevance This is a particularly useful citation for the project's $t_0/t$ scan. The paper often denotes the rung hopping by $t_\perp$; in this repository that symbol instead means interladder hopping. Keeping this notation difference explicit prevents the published ladder optimization from being mistaken for an interladder-coupling result.

\paper{Lin1998}{Lin, Balents, and Fisher 1998 --- Emergent SO(8) symmetry in a weakly interacting ladder}
\doi{10.1103/PhysRevB.58.1794}; \arxiv{cond-mat/9801285}
\summary A weak-coupling renormalization-group analysis shows that the half-filled two-leg ladder flows toward an integrable effective theory with an enlarged SO(8) symmetry. This description organizes the low-energy excitations, including spin, charge, and pairing-related modes, and gives detailed predictions for their gaps and quantum numbers. The authors also analyze lightly doped states and additional phases accessible with partly attractive interactions. The paper is a strong analytical demonstration that apparently different electronic responses can be closely related within the low-energy theory of a ladder, rather than being independent phenomenological channels.
\relevance It motivates comparing several sector-resolved gaps and treating charge, spin, and pairing as linked parts of the same microscopic problem. The emergent symmetry is derived in a specific weak-coupling limit, so neither its exact gap relations nor its degeneracies should be imposed on the project's $U/t=8$, finite-$V$, self-consistent array calculations.

\section{Isolated ladder benchmarks and the effects of nonlocal interactions}
These papers are the most useful comparators for bare-ladder measurements, the rung-hopping and $V$ scans, and interpretation of charge and pair profiles. They also expose the main finite-size and boundary ambiguities.

\paper{WhiteAffleckScalapino2002}{White, Affleck, and Scalapino 2002 --- Friedel oscillations and charge density waves}
\doi{10.1103/PhysRevB.65.165122}; \arxiv{cond-mat/0111320}
\summary This paper combines bosonization with DMRG analysis to distinguish boundary-induced density oscillations from true charge-density-wave order in chains and ladders. Open boundaries, which make DMRG efficient, can produce density patterns whose amplitude changes only weakly across the accessible system. The authors relate these patterns to correlation exponents, discuss their wave vectors, and develop an approach to extracting exponents from finite-size spectra with open boundaries. Their examples include Hubbard and $t$--$J$ chains and two-leg $t$--$J$ ladders, including a commensurate case with evidence for genuine charge order.
\relevance This is a central reference for evaluating the project's striped and charge-modulated states. A large Fourier peak or visually regular trough pattern is insufficient on its own; length dependence and bulk behavior matter. The paper's specific ordered $t$--$J$ example is not evidence that the present Hubbard point shares that phase, but its diagnostic reasoning applies directly.

\paper{Dolfi2015}{Dolfi and colleagues 2015 --- Pair correlations in doped Hubbard ladders}
\doi{10.1103/PhysRevB.92.195139}; \arxiv{1509.04709}
\summary High-accuracy DMRG revisits whether weakly doped repulsive two-leg Hubbard ladders are dominated by superconducting or density correlations. The calculations emphasize longer systems and careful extrapolation with truncation error, addressing earlier pair decays that were inconsistent with the expected low-energy theory. By comparing density and pair correlations with independently extracted Luttinger parameters, the authors recover the expected approach toward pairing-dominated behavior at very small doping. The studied isotropic ladder with $U/t=8$ makes this a close benchmark for the project's underlying building block. The central evidence is the asymptotic behavior of correlation functions, not a finite anomalous expectation value.
\relevance This is essential for the bare-ladder introduction and validation discussion. Its hopping and density conventions must be matched before importing numerical parameters. In particular, a Luttinger parameter defined through a density exponent need not equal a charge-mode parameter with a different normalization. The study does not determine the winner among this project's transverse mean-field branches.

\paper{Shen2023}{Shen, Zhang, and Qin 2023 --- Reexamining doped two-legged Hubbard ladders}
\doi{10.1103/PhysRevB.108.165113}; \arxiv{2303.16487}
\summary This DMRG study revisits ladder correlations using large systems, large retained-state spaces, and extrapolation with truncation error. A central finding is that the fitted superconducting exponent is highly sensitive to the reference bond's position under open boundaries. After controlling boundary and numerical effects, the authors find behavior consistent with a Luther--Emery liquid from very small hole doping through one-eighth doping, including reciprocal pairing and charge exponents in their convention. At larger dopings they report changes in these relationships and, eventually, qualitatively different ordering behavior. The paper provides a practical analysis of why plausible-looking correlation fits can nevertheless be misleading.
\relevance This is a high-priority benchmark for the project's density $n=0.9375$ and for moving reference bonds and fit windows. It supports using several observables together rather than accepting one exponent in isolation. Its bare-ladder conclusions still require rechecking when $t_0$, $V$, or self-consistent transverse fields differ from the studied model.

\paper{Zhou2023}{Zhou and colleagues 2023 --- Pairing robustness under intersite interactions}
\doi{10.1103/PhysRevB.108.195136}; \arxiv{2303.14723}
\summary DMRG calculations study an extended two-leg Hubbard ladder with next-nearest-neighbor hopping and density interactions extending to fourth neighbors. In the regimes examined, repulsive intersite interactions modify superconducting correlations quantitatively without immediately destroying their quasi-long-range character. An attractive nearest-neighbor term can substantially enhance pairing when its magnitude becomes comparable with the hopping, while sufficiently strong attraction leads to electron-hole phase separation. The analysis also examines charge and spin correlations, pairing structure, and the effects of onsite interaction. It therefore supplies a useful example in which the response to attraction is not monotonic improvement of a single superconducting phase.
\relevance This is directly relevant to the project's $V$ scan and to checking density inhomogeneity alongside pairing. The added diagonal hopping and longer-range interactions make its Hamiltonian different from the current ladder. Its boundaries and interaction thresholds should not be copied into the project phase diagram, and enhanced singlet correlations do not remove the need to check competing states.

\paper{Peng2023}{Peng and colleagues 2023 --- Enhanced pairing from near-neighbor attraction}
\doi{10.1103/PhysRevB.107.L201102}; \arxiv{2206.03486}
\summary This DMRG study examines an extended Hubbard model on long four-leg cylinders, motivated by experimental evidence for effective attraction in cuprate chains. Adding a near-neighbor attractive interaction strengthens long-distance superconducting correlations while weakening charge-density-wave correlations in the investigated parameter range. At sufficiently strong attraction within that range, pairing correlations become dominant and the associated superconducting susceptibility is strongly enhanced. The result provides a concrete numerical example of how an extra density interaction can shift the competition among electronic responses, connecting the experimental motivation for $V$ to an observable many-body effect beyond a single chain.
\relevance Cite it alongside Zhou and the ladder spectroscopy papers to motivate examining negative $V$. It is a four-leg cylinder study, so its correlation exponents and preferred charge pattern are not a direct prediction for a two-leg subsystem or a bulk mean-field array. Its title should not be paraphrased as proof that arbitrary attraction always stabilizes superconductivity.

\paper{Huang2025}{Huang and colleagues 2025 --- Pairing in an extended checkerboard Hubbard ladder}
\arxiv{2509.24415}\quad\textit{Preprint}
\summary This DMRG preprint investigates a checkerboard two-leg Hubbard ladder with different intra- and interplaquette hoppings and an attractive nearest-neighbor interaction within plaquettes. It reports that hopping inhomogeneity can reduce the attraction needed to reach a regime characterized by enhanced $d$-wave pairing correlations. The study examines the accompanying spin and single-particle behavior, linking changes in pairing to a broader reorganization of the low-energy state. Its results reinforce the idea that interaction strength and geometry should be considered jointly when trying to improve pairing in a ladder. The interpretation of superconductivity remains that appropriate to the finite-width, effectively one-dimensional system examined.
\relevance This is a recent comparison for combined hopping and $V$ tuning. Its principal density is $n=0.8125$, and its checkerboard modulation and intraplaquette attraction differ from the project's baseline and uniform interactions. The preprint's $t'$ denotes interplaquette hopping. Its correlation-based phase labels concern an effectively one-dimensional system, and the verified arXiv record is retained without an invented journal citation.

\paper{Kivelson2003}{Kivelson and colleagues 2003 --- How to detect fluctuating stripes}
\doi{10.1103/RevModPhys.75.1201}; \arxiv{cond-mat/0210683}
\summary This review develops ways to identify stripe correlations even when static long-range stripe order is absent. It discusses the interplay of fluctuating order, pinning, disorder, and experimental probes, explaining why local or induced spatial patterns can reveal a susceptibility without establishing a bulk broken symmetry. The broader lesson is that an ordering tendency is inferred through how the system responds and scales, rather than through visual inspection of a single pattern. The review also relates different charge and spin signatures, helping organize otherwise disparate observations in the cuprates.
\relevance Its reasoning is useful for the project's boundary-pinned textures, seeded branches, and bare-ladder response measurements. A density or spin profile can characterize a response to boundaries or fields even when it is not a spontaneous order parameter. Equal-time covariance is also distinct from static susceptibility: the latter contains excitation-energy weighting, so a large covariance eigenvalue is a candidate-selection signal, not a direct instability criterion.

\section{Coupled ladders and the transition from pairing to collective order}
These papers provide the conceptual link between a correlated one-dimensional building block and an ordered array. They distinguish pair transfer, single-particle coherence, and competition with other transverse instabilities.

\paper{Emery1997}{Emery, Kivelson, and Zachar 1997 --- The spin-gap proximity mechanism}
\doi{10.1103/PhysRevB.56.6120}; \arxiv{cond-mat/9610094}
\summary This theory treats conducting stripes coupled to a magnetically correlated environment. Pair hopping allows the mobile carriers to acquire a spin gap from nearby confined insulating regions, providing a route to a substantial pairing scale despite strong repulsion. Josephson coupling between stripes can then establish collective phase coherence at a separate scale. The work develops exact and controlled approximate results for a one-dimensional electron gas in an active environment and relates them to cuprate phenomenology. It is especially useful for distinguishing the local origin of pairs from the dimensional crossover that permits a superconducting array.
\relevance Cite it for the two-stage logic of pairing within a subsystem followed by coherence between subsystems. The project's arrays contain comparable Hubbard ladders, rather than precisely the stripe-plus-environment model studied here. The paper motivates effective pair transfer but does not derive the project's geometry-dependent kernels or show that density and magnetic competitors can be ignored.

\paper{KishineYonemitsu1998}{Kishine and Yonemitsu 1998 --- One-particle and two-particle dimensional crossovers}
\doi{10.1143/JPSJ.67.1714}; \arxiv{cond-mat/9802185}
\summary A perturbative renormalization-group treatment studies Hubbard ladders coupled by weak single-particle hopping. Growth of intraladder scattering toward a spin-gapped state suppresses coherent one-particle motion between ladders. In part of parameter space, a two-particle crossover therefore precedes the single-particle crossover and produces a superconducting instability. The authors compare the resulting flow and phase diagrams with coupled chains, showing why the low-energy structure of the building block matters for the dimensional crossover. This is an early explicit analysis of the hierarchy between transverse pair transfer and ordinary electronic band coherence.
\relevance This paper motivates the weak-interladder regime underlying a perturbative MPS+MF reduction. It also explains why small hopping relative to the bare leg bandwidth is not the only relevant comparison: the internal gaps matter. The calculation uses a perturbative treatment of intraladder interactions, so its numerical thresholds should not be imported into the project's strong-coupling ladder model.

\paper{GiamarchiTsvelik1999}{Giamarchi and Tsvelik 1999 --- Coupled ladders in a magnetic field}
\doi{10.1103/PhysRevB.59.11398}; \arxiv{cond-mat/9810219}
\summary Giamarchi and Tsvelik analyze spin ladders in a magnetic-field-induced gapless regime and then use their correlation functions to describe ordering in a weakly coupled three-dimensional array. The individual ladders retain one-dimensional quantum fluctuations, while the interladder coupling produces a collective transition. The analysis obtains a field-dependent transition temperature and related observables, including magnetic-resonance response. Although the microscopic degrees of freedom are spins rather than doped electrons, the construction is an important example of deriving higher-dimensional behavior from accurate low-energy information about a correlated one-dimensional subsystem.
\relevance It supplies conceptual support for combining nonperturbative intraladder physics with a transverse mean-field or susceptibility treatment. It is also relevant to the magnetic channel as an example of how interladder ordering can emerge. Its magnetic-field-driven spin model does not directly set the superconducting coupling, pair-binding denominator, or zero-field SDW boundary of the present fermionic array.

\paper{Arrigoni2004}{Arrigoni, Fradkin, and Kivelson 2004 --- Superconductivity in a striped Hubbard model}
\doi{10.1103/PhysRevB.69.214519}; \arxiv{cond-mat/0309572}
\summary This paper constructs a repulsive Hubbard model with strong spatial modulation that separates the system into nearly decoupled two-leg ladders with different hole densities. Using asymptotically controlled methods, the authors show how the ladder spin gap supplies a pairing scale and weak interladder coupling establishes phase order. The resulting ordering temperature depends on a positive power of the coupling between ladders. The construction gives a clear example of how a superconducting higher-dimensional state can be assembled from low-dimensional correlated components, while also showing why an appropriate pattern of inhomogeneity can be useful.
\relevance It is a direct conceptual predecessor for using ladder arrays to study repulsion-driven superconductivity. Its alternating dopings and explicitly modulated Hamiltonian differ from arrays of equivalent ladders. The project can use this paper to motivate the strategy, but must independently establish which order is energetically preferred when superconducting, charge, and magnetic channels are all retained.

\paper{Kantian2019}{Kantian and colleagues 2019 --- Repulsive pairing studied with parallel DMRG}
\doi{10.1103/PhysRevB.100.075138}; \arxiv{1903.12184}
\summary Massively parallel DMRG investigates anisotropic two-dimensional extended Hubbard systems relevant to quasi-one-dimensional superconductors. The authors use large strips and systematic extrapolation to examine the spin excitation spectrum, and combine those results with enhanced short-range correlations in a specific pairing channel. They build an indirect case for superconductivity while arguing against conventional magnetic order and a Fermi-liquid interpretation in the studied regime. The paper is valuable both as a numerical advance and as an example of assembling several distinct diagnostics when no single accessible observable settles the thermodynamic phase.
\relevance This is a useful comparison for the project's gap calculations and multichannel diagnosis. Its evidence for superconductivity is explicitly indirect, and the enhanced channel is $d_{xy}$ in that model rather than automatically the rung-versus-leg pattern of a two-leg Hubbard ladder. It does not implement the later fermionic MPS+MF construction and should be cited as complementary ancestry, not as the same algorithm.

\section{Matrix product states and numerical interpretation}
These references support the choice of solver and the interpretation of its finite-size and finite-entanglement results. Most belong in a methods section, with one or two citations in the introduction when motivating the computational approach.

\paper{White1992}{White 1992 --- Density matrix formulation for quantum renormalization groups}
\doi{10.1103/PhysRevLett.69.2863}
\summary White introduces the density-matrix renormalization-group idea: select a truncated subsystem basis using the reduced density matrix of a target state embedded in an environment. This changes how relevant states are retained compared with conventional real-space renormalization, producing highly accurate results for quantum chains. The method's success is demonstrated on Heisenberg systems, and it becomes the foundation for later fermionic ladder calculations. The central contribution is a principled truncation strategy tied to the entanglement of the target state, not an assumption that a small local basis can describe every correlated system equally well.
\relevance This is the original DMRG method citation for the paper. Its optimality statement concerns a particular truncation problem and should not be confused with guaranteed global convergence of a finite-bond-dimension calculation. It also supplies no error bound for the additional transverse mean-field approximation, which must be justified and assessed separately from the ladder solver.

\paper{Schollwock2011}{Schollw\"ock 2011 --- DMRG in the age of matrix product states}
\doi{10.1016/j.aop.2010.09.012}; \arxiv{1008.3477}
\summary This review explains DMRG in the language of matrix product states and develops the connections among canonical forms, local optimization, matrix product operators, symmetries, and truncation. It describes why the MPS representation is powerful for one-dimensional correlated systems and how its entanglement structure constrains accuracy. The review also covers extensions to dynamics and related algorithms, making it a general reference for interpreting the solver rather than just citing a software implementation. Its variational viewpoint helps distinguish the representation of a state from the procedure used to find that state.
\relevance Cite it alongside White when introducing the numerical building block. In this project, increasing the bond dimension tests the intraladder variational space, while improving self-consistency tests a different approximation step. Convergence of an energy or local field at one bond dimension does not automatically ensure converged long-distance correlations, correct branch selection, or accurate behavior of an infinite array.

\paper{CalabreseCardy2004}{Calabrese and Cardy 2004 --- Entanglement entropy and quantum field theory}
\doi{10.1088/1742-5468/2004/06/P06002}; \arxiv{hep-th/0405152}
\summary Calabrese and Cardy derive universal forms for the entanglement entropy of subsystems in one-dimensional quantum critical systems and related field theories. The logarithmic dependence on subsystem size is governed by the central charge, with different geometrical forms for infinite systems, finite systems, boundaries, and finite temperature. They also examine behavior away from criticality, where a finite correlation length changes the scaling. These results turn entanglement into a quantitative diagnostic of low-energy degrees of freedom rather than only a measure of numerical difficulty.
\relevance This is the theoretical citation for the project's rung-cut entropy fits and central-charge estimates. The boundary formula must match an open ladder and the interpretation requires a suitable scaling regime. Oscillations, finite length, and limited bond dimension can all bias a fitted coefficient. In addition, a self-consistent pairing field can gap modes of the effective ladder, so its fitted central charge need not equal the bare-ladder value.

\paper{Pollmann2009}{Pollmann and colleagues 2009 --- Finite-entanglement scaling at criticality}
\doi{10.1103/PhysRevLett.102.255701}; \arxiv{0812.2903}
\summary This work develops a quantitative theory of the artificial length scale introduced when a critical one-dimensional state is approximated by an MPS with finite bond dimension. The scaling depends on the central charge and the distribution of reduced-density-matrix eigenvalues. The authors test the theory at several critical points and explain why finite-entanglement scaling has a different origin from ordinary finite-size scaling. The result makes explicit that a numerically generated finite correlation length can arise from the restricted variational representation even when the underlying physical state is critical.
\relevance This is the key citation for distinguishing finite-$\chi$ effects from a physical gap or ordered phase. The project should vary length and bond dimension independently where conclusions depend on long-distance behavior. The paper's asymptotic scaling relations are not automatically quantitative for every open, spatially inhomogeneous, self-consistent ladder state, but the existence of a separate entanglement limitation is directly relevant.

\paper{Fishman2022}{Fishman, White, and Stoudenmire 2022 --- The ITensor software library}
\doi{10.21468/SciPostPhysCodeb.4}; \arxiv{2007.14822}
\summary This software paper describes ITensor's index-based interface for tensor-network calculations and its support for matrix product states, matrix product operators, tensor decompositions, and quantum-number-conserving block-sparse tensors. It explains how representing tensor connectivity explicitly reduces bookkeeping errors and enables implementations of common tensor-network algorithms. Examples connect the interface to many-body calculations, while the discussion of storage and symmetry support is relevant to practical performance and representation choices. The paper is the scholarly reference for the library family used by many modern DMRG implementations.
\relevance It belongs primarily in the methods or software acknowledgments of the eventual manuscript. Cite it together with the project's actual package versions and implementation details, because the publication does not describe every later ITensorMPS or CUDA feature. Using the library is not itself validation of the mean-field map, variational energy constants, or convergence criteria in a custom calculation.

\section{The closest MPS and mean field predecessors}
This is the narrowest and highest-priority group. The 2023 fermionic construction and 2025 multichannel extension directly motivate the project's method and scientific question; the 2026 preprint is a recent related application with a different ladder Hamiltonian.

\paper{Bollmark2020}{Bollmark, Laflorencie, and Kantian 2020 --- Dimensional crossover with MPS and mean field}
\doi{10.1103/PhysRevB.102.195145}; \arxiv{2005.02364}
\summary This study combines MPS calculations for one-dimensional subsystems with mean-field coupling between them to examine two- and three-dimensional arrays of hard-core and soft-core bosons. It investigates superfluid and insulating phases, including a first-order transition in the studied anisotropic setting, and calculates transition temperatures. Comparisons with analytical descriptions and quantum Monte Carlo for the full arrays provide benchmarks for the hybrid construction. A useful feature is its ability to remain practical where the available low-energy analytical treatment becomes less reliable. The paper establishes an important numerical foundation for the later fermionic MPS+MF developments.
\relevance Cite it when tracing the method's development and its benchmark history. Its bosonic tests should not be presented as a direct validation of every fermionic decoupling, frustrated transverse kernel, or competing SDW branch. The present work inherits the hybrid strategy, but its repulsive-ladder Hamiltonian and multichannel energy comparisons introduce additional model-specific questions.

\paper{Bollmark2023}{Bollmark and colleagues 2023 --- Fermionic MPS and mean field in two and three dimensions}
\doi{10.1103/PhysRevX.13.011039}; \arxiv{2207.03754}
\summary This paper develops fermionic MPS+MF by treating correlated one-dimensional subsystems numerically and weak transverse coupling through an effective mean-field construction. Attractive-Hubbard benchmarks are combined with field theory before the method is applied to three-dimensional arrays of repulsive two-leg Hubbard ladders. The ladder demonstration uses $U/t=8$ and $n=0.9375$, matching a central baseline of this project. It estimates superconducting transition temperatures using the relation between an effective excitation gap and the low-energy theory, with truncation and length extrapolations. Appendix E derives the ladder-array reduction and is the main methodological reference for the present implementation.
\relevance This is the closest starting point, but its superconductivity-focused ladder demonstration does not settle the project's full SDW/CDW/SC competition. The current zero-temperature variational energies are not transition temperatures or finite-temperature free energies. The paper's benchmark success also does not eliminate transverse mean-field error or justify applying the perturbative reduction outside the relevant gap hierarchy.

\paper{Bollmark2025}{Bollmark, K\"ohler, and Kantian 2025 --- Resolving charge and superconducting competition}
\doi{10.1103/PhysRevB.111.125141}; \arxiv{2301.08116}
\summary The authors extend fermionic MPS+MF to treat superconducting and charge-density channels together in three-dimensional arrays of attractive Hubbard chains with an additional nearest-neighbor interaction. They recover the expected symmetry-related SC/CDW coexistence at half filling and zero added interaction, then tune the balance with $V$ and doping. A crucial numerical feature is a CDW solution represented by alternating effective Hamiltonians and densities in a two-cycle of the self-consistency iteration. These alternating partners encode the charge arrangement of the array rather than real-time dynamics. The paper also discusses finite-size, bond-dimension, and boundary effects.
\relevance This is the direct predecessor for multichannel competition and the project's phase-resolved recurrence handling. Its benchmark uses negative-$U$ chains, not repulsive ladders with magnetic competitors. An arbitrary numerical two-cycle is therefore not automatically a physical CDW. The relevant contribution of the project is to test the expanded competition with its own geometry, spatial diagnostics, and consistent energetic comparisons.

\paper{Kohler2026}{K\"ohler and Kantian 2026 --- Predicting ordering scales in quantum simulators}
\arxiv{2608.18861}\quad\textit{Preprint submitted 19 August 2026}
\summary This recent preprint applies MPS+MF and low-energy field theory to arrays of mixed-dimensional $t$--$J$ ladders designed for cold-atom quantum simulators. It characterizes pair binding, spin gaps, charge response, and finite-size effects of the isolated building blocks before estimating superconducting ordering scales for the coupled system. A central result is that larger binding energy need not maximize the estimated transition temperature: spin gaps, stiffness, susceptibility, and pair mobility can move in competing directions. The work also identifies finite systems for which the assumptions entering the low-energy description become unreliable.
\relevance This is a timely comparison for the project's fixed-sector backbone and perturbative gap checks. Its suppressed rung tunneling and exchange-engineered ladder differ from the current Hubbard Hamiltonian, and its temperature estimates remain within the stated theoretical treatment. Cite it as a preprint and avoid treating it as evidence that larger $|E_p|$, stronger pairing fields, or a chosen finite length automatically improve bulk superconductivity.

\section{Using the review in the introduction}
The introduction can be organized into the following sequence. The aim is to explain the scientific question and its relation to prior work without turning the introduction into a numerical audit.

\begin{enumerate}
\item \textbf{Start with the correlated-electron problem.} Introduce repulsion-driven pairing and the competing charge and magnetic responses of doped Mott systems. A compact citation set is Hubbard, Scalapino, Arovas, and Fradkin \citep{Hubbard1963,Scalapino2012,Arovas2022,Fradkin2015}. Anderson and Lee provide optional historical depth.
\item \textbf{Explain why the ordering question remains difficult.} Use the two-dimensional numerical work to motivate sensitivity to Hamiltonian parameters and low-energy textures \citep{Zheng2017,Qin2020,Xu2024}. State model differences explicitly rather than presenting these papers as contradictory answers to one identical calculation.
\item \textbf{Introduce ladders as physical building blocks.} Connect the spin-gapped ladder picture to the original pairing results and modern correlation benchmarks \citep{DagottoRice1996,BalentsFisher1996,Dolfi2015,Shen2023}. Uehara and Nagata establish the material and dimensional-crossover motivation \citep{Uehara1996,Nagata1998}.
\item \textbf{Motivate the tuning parameters.} Rung hopping changes ladder electronic structure, while effective nonlocal attraction has both spectroscopic and numerical motivation \citep{Noack1997,Chen2021,Padma2025,Zhou2023,Peng2023}. Keep material inference, isolated-ladder response, and array ordering separate.
\item \textbf{Present the coupled-subsystem strategy.} Explain why treating the ladder accurately and its weak transverse coupling approximately is useful \citep{Arrigoni2004,Bollmark2020,Bollmark2023}. Introduce multichannel competition with the 2025 extension \citep{Bollmark2025}. The 2026 preprint supplies a recent connection to gap hierarchies and simulator design \citep{Kohler2026}.
\item \textbf{State this project's question precisely.} The natural question is how transverse connectivity, rung hopping, and nearest-neighbor interactions alter the competition among paired, charge-modulated, and magnetic solutions of repulsive ladder arrays. Describe any answer in terms of the calculations actually completed. Do not claim a thermodynamic phase diagram, a critical temperature, or a first-ever result solely from this review.
\end{enumerate}

\subsection*{Position relative to the nearest work}
The literature supports a focused positioning: extend the repulsive-ladder MPS+MF setting to geometry-dependent competition among superconducting, charge, and magnetic textures, with comparable zero-temperature energies and spatial diagnostics. This combines the ladder setting of Ref.~\citep{Bollmark2023} with the multichannel motivation of Ref.~\citep{Bollmark2025}; it is not a claim that the MPS+MF framework itself is new. The isolated-ladder and extended-interaction papers supply the reasons to vary $t_0$ and $V$. A final novelty statement should be checked against the completed results and the literature available at submission.

\subsection*{Project evidence update: 19 September 2026}

The synchronized project data now include the full unfrustrated-cubic grid,
six finer square coordinates, all four repulsive-$V$ square starts
and all four trellis starts. These are local results, not findings
of the papers annotated above; this update adds no literature search.
Both seeds reach stripes at all nine cubic coordinates, including
$(t_0/t,V/t)=(1.4,-0.4)$ and $(1.4,-0.2)$, where square remains paired.
At square $(1.2,+0.2)$ all four starts lose anomalous pairing, including
both intertwined wavelengths. The period-eight seed leaves six unevenly
spaced spin nodes and a higher diagnostic endpoint energy, consistent with
incomplete coarsening or magnetic defects. Repulsive $V$ therefore cannot
be presented as a general mechanism for intertwining based on the legacy
frustrated-cubic result. Both one-ladder trellis seeds reach opposite-sign
leg/rung pairing at $(1,0)$, but both rectangular two-ladder seeds develop
essentially unpaired stripes. Large alternating field/profile changes remain
in the latter, despite nearly constant spin RMS. This is numerical relaxation,
not physical dynamics or an accepted orbit. Different spatial constraints
and open-end cuts qualify the cell comparison. Geometry, spatial ansatz and
competing channels all matter when relating ladder tendencies to array order.

The finer square cuts retain distinct textures at $(1.25,-0.4)$ and
$(1.4,-0.05)$. Their full variational energies are nearly coincident on the
absolute scale. Along $t_0=1.4$ a small slope downturn appears in the final
$-0.05\le V/t\le0$ interval; the $V=-0.4$ cut is smoothly curved in $t_0$
at the present resolution. Neither demonstrates a discontinuous derivative
or an accepted energy crossing. Weak spin in the paired $V=-0.05$ run is
strongly concentrated near the open ends, so bulk coexistence is not
established. Matched physical spin/pairing cuts now accompany the energy
curves in the manuscript. All 38 completed run histories remain unaccepted
at $L=64$ and $\chi=200$. The detailed energy curves, spatial profiles and
provenance are in \path{reports/campaign_review_20260918/README.md}; the
updated manuscript develops their interpretation. These qualifications
should accompany comparisons with the intertwined-order literature and
Refs.~\citep{Bollmark2023,Bollmark2025}.

\subsection*{Conventions that need care when drafting}
\begin{itemize}
\item \textbf{Hopping symbols.} Many ladder papers call the rung hopping $t_\perp$. Here it is $t_0$, and $t_\perp$ means hopping between ladders. A symbol $t'$ may instead denote diagonal or interplaquette hopping in another paper.
\item \textbf{Binding and gap conventions.} The repository's hole-binding energy is negative for binding, while several MPS+MF papers define a positive pair-breaking scale. Match particle sectors and signs before comparing numbers. Pair binding, spin excitation gaps, and a charge gap are different quantities.
\item \textbf{Luttinger parameters.} State the operator normalization and the powers of the actual correlation functions. Do not equate a density exponent, a charge-mode parameter, and a stiffness-derived parameter merely because each is called $K_\rho$ or $K_c$.
\item \textbf{Order and self-consistency.} Pair binding, algebraic isolated-ladder pairing, a self-consistent anomalous field, and thermodynamic superconducting order are different statements. A two-cycle in an iteration is likewise distinct from time-dependent physics and requires a spatial interpretation.
\item \textbf{Comparison scope.} Different transverse geometries are different Hamiltonians. Competing branches can be ranked at one Hamiltonian point using a common energy functional; energies from different geometries should not be described as competing phases of a single unchanged model.
\end{itemize}

\bibliographystyle{unsrtnat}
\begingroup
\raggedright
\renewcommand{\bibpreamble}{\interlinepenalty=10000}
% Support both project-root (Overleaf) and literature-directory local builds.
\IfFileExists{literature/references.bib}{%
  \bibliography{literature/references}%
}{%
  \bibliography{references}%
}
\endgroup
\end{document}
